\section{Prescribing a positive thickness}\label{sec:sheets-boxes}
We shall also need boxes of a prescribed positive thickness, for which
an unbounded choice of layer heights is convenient.
\begin{lemma}[Exponential scaling of full boxes]
\label{lem:sa-full-box-layers}\label{lem:sb-boxes}
Suppose $R_i^\pm\subset\mathbb R^2$, $1\le i\le N$, $N\ge1$,
are separately separated nondegenerate rational closed rectangles.
Let $F_i$ be their positive diagonal affine maps, and explicitly require
$F_i(R_i^-)=R_i^+$. Their factors $\lambda_{i1},\lambda_{i2}$ are
positive rational numbers, given by the corresponding rational side-length
ratios. Put
$\lambda_{i3}=(\lambda_{i1}\lambda_{i2})^{-1}$, and choose any positive
rational source half-heights $h_i^-$. Define $h_i^+=\lambda_{i3}h_i^-$.
There is an effective compactly supported divergence-free unit-time pulse
whose map on each full box $R_i^-\times[-h_i^-,h_i^-]$ is the specified
horizontal map together with $z\mapsto\lambda_{i3}z$. It maps that box
to $R_i^+\times[-h_i^+,h_i^+]$, and the formula is exact on a positive
three-dimensional neighborhood. Its centers interpolate linearly in the
horizontal coordinates, while every half-side is at most the larger of
its endpoint half-sides. The pulse vanishes near both time endpoints.
\end{lemma}
\begin{proof}
Let $Z=\max_{i,\pm}h_i^\pm$. Choose a positive rational collar
$0<\eta\le1/4$ no larger than one quarter of each within-family horizontal
coordinate-separating gap. If a family has only one member, there is no
pairwise condition. Assign heights $z_i=4(Z+1)i$. More generally, the
same proof permits any heights whose pairwise gaps exceed $2Z+2\eta$.
With successive flat switches $s_1,s_2,s_3$, lift, deform, and lower using
\[
 C_i=((1-s_2)c_i^-+s_2c_i^+,z_i(s_1-s_3)),\qquad
 D_i=\diag(\lambda_{i1}^{s_2},\lambda_{i2}^{s_2},\lambda_{i3}^{s_2}).
\]
The middle-stage vertical half-size is at most $Z$, because every
positive geometric interpolation stays between its two endpoint values.
During ascent and descent, the fixed source and target projections,
respectively, keep the $\eta$ enlargements disjoint. During deformation
the height condition separates those enlargements. Product cutoffs equal
to one on the $\eta/2$ enlargement and supported in the $\eta$
enlargement therefore meet Lemma~\ref{lem:sa-curl-motion}.
The determinant of $D_i$ is one. That lemma gives the exact affine path
on every full box and a positive initial neighborhood. Only finitely
many bounded paths occur, so the support lies in a fixed compact set.
Logarithms and exponentials of positive rational factors are effective.

This proof permits different initial thicknesses because it uses only
the finite maximum of all endpoint half-heights. In the special case
$h_i^-=1$ one has $Z=\max_i(1,\lambda_{i3})$ and the stated default
heights are $4(Z+1)i$. Unlike the linear-width sheet path, this
interpolation does not make each individual horizontal trajectory a
convex combination of its endpoints; the center and half-width bounds
are the assertions used below.
\end{proof}

\subsection{Records attached to individual work cells}\label{subsec:sa-cell-records}
A related recorder keeps a word of history immediately after each work
cell.  This changes the actual rule cylinders: a stay writes its record
on the right, and a right move transfers its newly written history in a
separate scan.  It is not a renaming of \eqref{eq:sa-interleaved-table}.

Use an origin-tagged work alphabet $\Sigma$, moves $D=\{0,+,-\}$, and
an incoming tag $e\in D\sqcup\{*\}$.  Extend halt states by idle
instructions.  The controls are $C(q,e),F(p),G(p)$; the distinct records
are $\gamma(q,e,a)$, all outside $\Sigma$; write $\mathcal G$ for
their alphabet.  Initialize at $C(q_0,*)$ with stacks $\square^\infty$ and
$w\square^\infty$, where $w$ includes the permanent origin tag.
At a main control with head at $j$, its invariant is
\[
 L=\operatorname{rev}(a_{j-1}W_{j-1})
   \operatorname{rev}(a_{j-2}W_{j-2})\cdots,
 \qquad R=a_jW_j\,a_{j+1}W_{j+1}\cdots,
\]
where every $W_i$ is a finite record word, initially empty.  For an instruction writing $b$ and retaining its record $g$, the main
update changes $a_jW_j$ to $bgW_j$.  The $F$ scan transfers the whole record word $gW_j$
before exposing the next work cell; the $G$ scan transfers the previous
cell's reversed record and then its work symbol.  Thus every instruction
finishes in finitely many positive steps and gives the stated invariant.
Erasing records and using the origin tag recovers the original tape.

If $\delta(q,a)=(p,b,d)$ and $g=\gamma(q,e,a)$,
use
\begin{equation}\label{eq:sa-cell-records-3}
\begin{array}{ll}
 d=0:&(C(q,e);\epsilon,a)\to(C(p,0);\epsilon,bg),\\
 d=+:&(C(q,e);\epsilon,a)\to(F(p);b,g),\\
 d=-:&(C(q,e);\epsilon,a)\to(G(p);\epsilon,bg).
\end{array}
\end{equation}
For each record $g'$ and work symbol $a'$, add
\begin{equation}\label{eq:sa-cell-records-4}
\begin{array}{ll}
 (F(p);\epsilon,g')\to(F(p);g',\epsilon),&
 (F(p);\epsilon,a')\to(C(p,+);\epsilon,a'),\\
 (G(p);g',\epsilon)\to(G(p);\epsilon,g'),&
 (G(p);a',\epsilon)\to(C(p,-);\epsilon,a').
\end{array}
\end{equation}


The sources are separated by their tested top symbols.  Image cylinders
into $C(p,0)$ are separated by their right prefixes $bg$, and those into
$C(p,+)$ or $C(p,-)$ by the exiting work symbol.  No rule enters $C(p,*)$.
Into $F(p)$ a main entry has work first on $L$ and its identifying record
first on $R$, whereas a loop has a record first on $L$.  Into $G(p)$ a
main entry has $bg$ first on $R$, whereas a loop has a record first there.
These distinctions recover the incoming rule for all tails.

\begin{theorem}[Input shifts and periodic full-box motion]
\label{thm:sa-cell-records}
This recorder gives an effective smooth general force on $\mathbb R^3$
with fixed compact spatial support, period one from $t=0$, and all mixed
derivatives bounded.  Its zero-data solution is globally smooth with
uniformly bounded energy.  The particle starting at the origin has first
coordinate greater than $2$ at some time exactly when the machine halts.
Uniqueness holds for smooth velocities in $C_tH^2\cap C_t^1L^2$ with
velocity and gradient bounded on finite intervals, and pressure in $H^1$
modulo constants at each time.  Pressure is determined modulo a function
of time.
\end{theorem}
\begin{proof}
Take odd digits and base $B=2|\Sigma\sqcup\mathcal G|+1$, where
$\mathcal G$ is the record alphabet.  Put
$l_0=E(\square^\infty)$ and $r_0=E(w\square^\infty)$.
Give $C(q_0,*)$ offset zero, other nontrigger controls distinct offsets
$-4,-8,\ldots$, and trigger controls
$C(q,d)$ with $q$ halting and $d\in D$ offsets $4,8,\ldots$.
Encode by
\begin{equation}\label{eq:sa-cell-records-5}
 (o_s-l_0+E(L),-r_0+E(R),0).
\end{equation}
The initial code is exactly zero.  For every rule use its two prefix
rectangles, with scales
$\lambda_1=B^{|v|-|v'|}$, $\lambda_2=B^{|w|-|w'|}$ and
$\lambda_3=(\lambda_1\lambda_2)^{-1}$.
Lemma~\ref{lem:sa-full-box-layers} acts on the entire box
$R^-\times[-1,1]$, sending it to $R^+\times[-\lambda_3,\lambda_3]$.
The encoded central sheet returns to the same plane $z=0$ on successive
periods; the full box thickness generally changes.  One may take margin $B^{-m_*}/10$,
where $m_*$ is the maximum prefix length, and heights $4(Z+1)i$ as in
that lemma.  Periodize the resulting pulse from time zero; the time
collars give $U(0)=0$.  Lemma~\ref{lem:sa-realization} and its pointwise
$H^1$ pressure refinement give the equation and uniqueness.

The code induction now begins at $t=0$.  A reached halt has a trigger
code with first coordinate at least $4-l_0>2$.  An initial halt reaches
one after its first idle instruction, changing tag $*$ to $0$.
For a nonhalting run, every used source and target has first-coordinate
center at most $1$ and half-width at most $1/2$.  The exponential box
path has convexly interpolated centers and half-width at most the larger
endpoint half-width; hence it stays below $3/2<2$ throughout.  Input
shifts affect the finite spatial prescription, and no loader is required.
\end{proof}

